\subsection{Sprint 4}

Na sprint final, continuei a ajustar a tela de anexo e a integrar os
componentes necessários. Também tentei desenvolver a lógica de download de
arquivos do componente de documentos. Essa funcionalidade se mostrou
trabalhosa e complexa demais, consumindo muito tempo com pouco progresso. Por
isso, minha dupla foi realocada para refatorar a tela inicial do aplicativo.

Depois que a tela inicial ficou pronta, fui designado para a parte Web do
projeto. Ajudei no desenvolvimento da tela de listagem dos procedimentos
criminais, novamente com o colega Gustavo Pretto. Passei a maior parte desse
tempo corrigindo a integração com as rotas do \textit{backend} e os erros no
próprio ambiente de desenvolvimento.
\\
\\
\fbox{%
    \parbox{1\linewidth}{%
        \setlength{\parindent}{15pt}
            \itshape A última sprint foi relativamente tranquila em comparação
com o que havia acontecido antes no projeto. A atitude do time já tinha mudado
e senti que a maioria conseguia trabalhar ou pedir ajuda quando necessário.
Nos dois últimos fins de semana, o time cooperou para resolver alguns detalhes
finais.

Se eu pudesse recomeçar este projeto, aproveitaria melhor a experiência:
poderia ter tomado a iniciativa antes e me comunicado mais. Apesar de a
primeira metade ter sido desafiadora e intimidadora, acho que ela me trouxe
lições importantes para minhas futuras AGES.
    }%
}
